\documentclass[11pt,a4paper]{article}

% -- Encoding & Fonts --
\usepackage[utf8]{inputenc}

% -- Mathematics --
\usepackage{amsmath,amssymb,amsthm,mathtools}

% -- Layout --
\usepackage[margin=2.5cm]{geometry}
\usepackage{setspace}
\doublespacing
\usepackage{parskip}

% -- Line numbers for review --
\IfFileExists{lineno.sty}{%
  \usepackage{lineno}
  \linenumbers
}{}

% -- Section formatting --
\usepackage{titlesec}
\titleformat{\section}{\large\bfseries}{\thesection.}{0.5em}{}
\titleformat{\subsection}{\normalsize\bfseries}{\thesubsection}{0.5em}{}
\titleformat{\subsubsection}{\normalsize\itshape}{\thesubsubsection}{0.5em}{}

% -- TOC formatting --
\usepackage[titles]{tocloft}

% -- Headers & footers --
\usepackage{fancyhdr}
\pagestyle{fancy}
\fancyhf{}
\fancyhead[L]{\small\textit{GCC: A dynamical regime model for global access}}
\fancyhead[R]{\small\textit{Tugan}}
\fancyfoot[C]{\thepage}
\renewcommand{\headrulewidth}{0.4pt}

% -- Hyperlinks --
\usepackage{xcolor}
\usepackage[colorlinks=true,linkcolor=black,citecolor=black,urlcolor=black]{hyperref}

% -- Bibliography --
\usepackage[numbers,sort&compress]{natbib}

\usepackage{graphicx}

% -- Theorem environments --
\theoremstyle{plain}
\newtheorem{theorem}{Theorem}[section]
\newtheorem{proposition}[theorem]{Proposition}
\newtheorem{corollary}[theorem]{Corollary}
\theoremstyle{definition}
\newtheorem{definition}[theorem]{Definition}
\theoremstyle{remark}
\newtheorem*{remark}{Remark}
\newtheorem*{significance}{Significance}

% -- Custom commands --
\newcommand{\Deff}{D_{\mathrm{eff}}}
\newcommand{\Mtau}{M_{\tau}}
\newcommand{\Rmin}{R_{\min}}
\newcommand{\Dmin}{D_{\min}}
\newcommand{\Dmax}{D_{\max}}
\newcommand{\Mmin}{M_{\min}}
\newcommand{\Mmax}{M_{\max}}
\newcommand{\Gmax}{G_{\max}}
\newcommand{\Kcop}{K_c^{\mathrm{op}}}
\newcommand{\Lobs}{L_{\mathrm{obs}}}

% -- Title block --
\title{%
  \textbf{The Gamma-Code Consensus: a dynamical regime model for level of consciousness and global access}
}
\author{%
  Mehmet Oktay Tugan\textsuperscript{1,*}
}
\date{}

% ==================================================================
\title{\textbf{Supplementary Material for: The Gamma-Code Consensus}}
\author{Mehmet Oktay Tugan}
\date{}
\begin{document}
\maketitle
\noindent This supplementary file contains the mathematical appendix and code-resource description moved out of the main manuscript for the Neuroscience of Consciousness submission version.

\appendix
\makeatletter
% Inline implementation of removefromreset (avoids needing the remreset package):
% remove section from the reset-list of the theorem counter so numbering is globally unique
\providecommand*{\@removefromreset}[2]{%
  \let\@tempa\@elt
  \let\@elt\relax
  \expandafter\let\expandafter\@tempb\csname cl@#2\endcsname
  \expandafter\def\csname cl@#2\endcsname{}%
  \def\@elt##1{%
    \ifx ##1#1\else
      \expandafter\@temptokena\expandafter{\csname cl@#2\endcsname}%
      \edef\x{\noexpand\expandafter\noexpand\def\noexpand\csname cl@#2\noexpand\endcsname{\the\@temptokena \noexpand\@elt{##1}}}\x
    \fi}%
  \@tempb
  \let\@elt\@tempa}
\@removefromreset{theorem}{section}
\makeatother
\section*{Mathematical Appendix}
\addcontentsline{toc}{section}{Mathematical Appendix}
\label{sec:appendix}
\setcounter{section}{0}
\setcounter{theorem}{0}
\renewcommand{\thesection}{A.\arabic{section}}
\renewcommand{\thesubsection}{A.\arabic{section}.\arabic{subsection}}
\renewcommand{\thetheorem}{A.\arabic{theorem}}

% Redefine section format for appendix: no trailing dot (already in thesection)
\titleformat{\section}{\large\bfseries}{\thesection}{0.6em}{}
\titlespacing*{\section}{0pt}{1.5em}{0.8em}

% Widen TOC number column for appendix entries (A.1 needs more space than 1)
\addtocontents{toc}{\string\setlength{\string\cftsecnumwidth}{3.2em}}

This appendix supplements the GCC with mathematical results that address three common lines of attack: (1)~Is the dynamics well-posed? (2)~Are the observables controlled? (3)~Can the gain-based re-entry idea be formulated as a sufficient condition?

\paragraph{Core definitions used below.}
For self-contained reading, the core dynamics, ridge regularization, and residual-backbone definition from the main manuscript are repeated here:
\begin{equation}
\label{eq:dynamics}
\mathrm{d}\theta_i(t)
= \biggl[\,\omega_i
+ K \cdot G(t) \cdot g_i(t) \sum_{j} W_{ij}(t)\,\sin\!\bigl(\theta_j(t) - \theta_i(t) - \varphi_{ij}(t)\bigr)
\biggr]\,\mathrm{d}t
+ \sigma_i(t)\,\mathrm{d}B_i(t).
\end{equation}
\begin{equation}
\label{eq:ridge}
\varepsilon_D^{\lambda}(t) =
\max\!\biggl(\lambda\,\frac{\mathrm{tr}\bigl(\hat{C}_D(t)\bigr)}{N},\;\varepsilon_{\mathrm{floor}}\biggr),
\qquad
C_D^{\lambda}(t) = \hat{C}_D(t) + \varepsilon_D^{\lambda}(t)\cdot I.
\end{equation}
\begin{equation}
\label{eq:backbone}
S(t;\,g_*) = \text{largest connected component of the subgraph on } \{i : g_i(t) \geq g_*\}.
\end{equation}


\section{Assumptions}
\label{app:voraussetzungen}

We consider, over a finite horizon $t \in [0,T]$, the GCC dynamics~\eqref{eq:dynamics} under the following assumptions.

\paragraph{A1.\;Structure.}
$W_{ij}(t) \geq 0$, \;\;$\sum_j W_{ij}(t) = 1$ for all $i$ with non-empty rows, \;\;$W_{ii}(t) = 0$.

\paragraph{A2.\;Boundedness.}
There exist finite constants $\Gmax$, $\Sigma_{\max}$, $\Phi_{\max}$ such that
\[
0 \leq G(t) \leq \Gmax, \quad
0 \leq g_i(t) \leq 1, \quad
0 \leq \sigma_i(t) \leq \Sigma_{\max}, \quad
|\varphi_{ij}(t)| \leq \Phi_{\max}.
\]

\paragraph{A3.\;Temporal regularity.}
All control functions are measurable and bounded; for stability statements, piecewise continuity suffices.

\paragraph{A4.\;Observation windows.}
$\Deff^{\lambda}$ and $\Mtau$ are computed on fixed window lengths $\tau_D > 0$ and $\tau_M > 0$, with ridge regularization according to~\eqref{eq:ridge}.


\section{Elementary Bounds of the Observables}

\begin{proposition}[Elementary bounds]
\label{prop:schranken}
For all times $t$:
\begin{equation}
0 \leq R(t) \leq 1, \qquad
1 \leq \Deff(t) \leq N, \qquad
0 \leq \Mtau(t) \leq \tfrac{1}{4}.
\end{equation}
\end{proposition}

\begin{proof}[Proof sketch]
$R(t)$ is the modulus of a mean of unit complex numbers. $\Deff(t)$ is the participation ratio of a positive semi-definite matrix. $\Mtau(t)$ is the variance of a time series with values in $[0,1]$; its maximum is $\frac{1}{4}$.
\end{proof}

\begin{significance}
All three core observables are mathematically compact and controlled. This prevents empirical fits from drifting into unphysical regions.
\end{significance}


\section{Well-Posedness of the Dynamics}

\begin{theorem}[Global Lipschitz property and well-posedness]
\label{satz:lipschitz}
Define the drift vector $b(\theta,t)$ component-wise by
\begin{equation}
b_i(\theta,t) = \omega_i + K\,G(t)\,g_i(t)\sum_j W_{ij}(t)\,\sin\!\bigl(\theta_j - \theta_i - \varphi_{ij}(t)\bigr).
\end{equation}
Then, for all $\theta,\psi \in \mathbb{R}^N$ and all $t \in [0,T]$,
\begin{equation}
\label{eq:lipschitz}
\|b(\theta,t) - b(\psi,t)\|_{\infty} \;\leq\; 2\,|K|\,\Gmax\,\|\theta - \psi\|_{\infty}.
\end{equation}
Moreover, $\|b(\theta,t)\|_{\infty} \leq \|\omega\|_{\infty} + |K|\,\Gmax$. In particular, under \textup{A1--A3}, the GCC dynamics has a unique global strong solution on $[0,T]$ for every initial condition $\theta(0)$.
\end{theorem}

\begin{proof}[Proof sketch]
For each component, using $|\sin(x) - \sin(y)| \leq |x - y|$:
\begin{align}
|b_i(\theta,t) - b_i(\psi,t)|
&\leq |K|\,\Gmax \sum_j W_{ij}(t)\,\bigl|(\theta_j - \theta_i) - (\psi_j - \psi_i)\bigr| \notag\\
&\leq |K|\,\Gmax\,\Bigl[\sum_j W_{ij}(t)\,|\theta_j - \psi_j| + |\theta_i - \psi_i|\underbrace{\sum_j W_{ij}(t)}_{=\,1}\Bigr] \notag\\
&\leq 2\,|K|\,\Gmax\,\|\theta - \psi\|_{\infty}. \label{eq:lip_detail}
\end{align}
Standard results for SDEs with globally Lipschitz drift and bounded diffusion yield existence and uniqueness.
\end{proof}


\section{Finite-Horizon Stability}

\begin{theorem}[Stability under drift perturbations]
\label{satz:stabilitaet}
Consider two systems driven by the same Brownian motion:
\begin{align}
\mathrm{d}\theta(t) &= b\bigl(\theta(t),t\bigr)\,\mathrm{d}t + \Sigma(t)\,\mathrm{d}B(t), \\
\mathrm{d}\vartheta(t) &= \tilde{b}\bigl(\vartheta(t),t\bigr)\,\mathrm{d}t + \Sigma(t)\,\mathrm{d}B(t).
\end{align}
Suppose both drifts are globally $L$-Lipschitz and
\begin{equation}
\|b(x,t) - \tilde{b}(x,t)\|_{\infty} \leq \varepsilon(t) \quad \text{for all } x \in \mathbb{R}^N,\; t \in [0,T].
\end{equation}
Then
\begin{equation}
\label{eq:gronwall}
\boxed{\;
\|\theta(t) - \vartheta(t)\|_{\infty}
\;\leq\; e^{Lt}\,\biggl(\|\theta(0) - \vartheta(0)\|_{\infty} + \int_0^t \varepsilon(s)\,\mathrm{d}s\biggr)
\;}
\end{equation}
\end{theorem}

\begin{proof}[Proof sketch]
The diffusion terms cancel in the difference equation. For $\delta(t) = \theta(t) - \vartheta(t)$, we have $\|\delta'(t)\|_{\infty} \leq L\,\|\delta(t)\|_{\infty} + \varepsilon(t)$. Gronwall's inequality yields the claim.
\end{proof}

\begin{significance}
Small changes in $W$, $G$, $g_i$, or $\varphi_{ij}$ produce only controlled changes in the trajectory on finite time horizons.
\end{significance}


\section{Continuity of the Observables}

\begin{proposition}[Local Lipschitz property]
\label{prop:stetigkeit}
For fixed $\tau_D$, $\tau_M$, $\lambda$, and $\varepsilon_{\mathrm{floor}}$, the observables $R(t)$, $\Deff^{\lambda}(t)$, and $\Mtau(t)$ are continuous functions of the window trajectory. On compact sets they are locally Lipschitz. In particular, on any suitable compact tube there exists a constant $\Lobs$ such that
\begin{equation}
\label{eq:Lobs}
\|O(\mathrm{traj}_1) - O(\mathrm{traj}_2)\|_{\infty} \;\leq\; \Lobs \cdot \|\mathrm{traj}_1 - \mathrm{traj}_2\|_{\infty},
\end{equation}
where $O = (R,\;\Deff^{\lambda},\;\Mtau)$.
\end{proposition}

\paragraph{Empirical estimation of $\Lobs$.}
One chooses small initial perturbations $\delta_0$, simulates $\theta$ and $\theta^{\delta}$ with the same noise path, and computes
\begin{equation}
\hat{L}_{\mathrm{obs}}(\delta_0) = \frac{\sup_t \|O(\theta^{\delta})(t) - O(\theta)(t)\|_{\infty}}{\sup_t \|\theta^{\delta}(t) - \theta(t)\|_{\infty}}.
\end{equation}


\section{Residual-Backbone Re-Entry}

\begin{theorem}[Sufficient condition for re-entry]
\label{satz:reentry}
Let $S = S(t_1;\,g_*)$ be the residual backbone structure as in~\eqref{eq:backbone}. Consider, on $I = [t_1,t_2]$, the system restricted to $S$ in which the coupling terms coming from outside $S$ are switched off. Define the external residual mass
\begin{equation}
\alpha_S(t) = \max_{i \in S}\sum_{j \notin S} W_{ij}(t).
\end{equation}
Assume that:
\begin{enumerate}
\item The reduced system lies on $I$ within the access region with safety margin $\rho > 0$:
\begin{align}
R_S(t) &> \Rmin + \rho, \notag\\
\Dmin + \rho &< D_{\mathrm{eff},S}^{\lambda}(t) < \Dmax - \rho, \label{eq:marge}\\
\Mmin + \rho &< M_{\tau,S}(t) < \Mmax - \rho. \notag
\end{align}
\item The initial states coincide at $t_1$.
\item A local Lipschitz constant $\Lobs$ for $O = (R_S,\,D_{\mathrm{eff},S}^{\lambda},\,M_{\tau,S})$ exists.
\item The following holds:
\begin{equation}
\label{eq:eta}
\boxed{\;
\int_{t_1}^{t_2} \alpha_S(s)\,\mathrm{d}s \;<\; \eta, \qquad
\eta \;<\; \frac{\rho}{\Lobs \cdot |K| \cdot \Gmax \cdot e^{2|K|\Gmax(t_2 - t_1)}}
\;}
\end{equation}
\end{enumerate}
Then the full system restricted to $S$ remains in the access region throughout $I$.
\end{theorem}

\begin{proof}[Proof sketch]
By Theorem~\ref{satz:stabilitaet} with $L = 2|K|\Gmax$ and $\varepsilon(t) = |K|\Gmax\,\alpha_S(t)$,
\begin{equation}
\sup_{t \in I}\|\theta_{\mathrm{full}}(t) - \theta_{\mathrm{red}}(t)\|_{\infty}
\;\leq\; |K|\,\Gmax\,e^{2|K|\Gmax(t_2 - t_1)}\cdot\eta.
\end{equation}
Using $\Lobs$, it follows that $\|O_{\mathrm{full}} - O_{\mathrm{red}}\|_{\infty} \leq \Lobs \cdot |K|\,\Gmax\,e^{2|K|\Gmax(t_2-t_1)}\cdot\eta$. If $\eta$ is chosen according to~\eqref{eq:eta}, this stays below $\rho$.
\end{proof}

\begin{significance}
This is the mathematically sharpest form of the gain-based terminal-lucidity hypothesis: a transient re-entry is sufficiently possible whenever a residual subnetwork can carry the dynamics and the rest of the system does not perturb it too strongly on a short interval.
\end{significance}

\paragraph{Conservativity.}
The $\eta$-condition is a Gronwall-based worst-case bound. If numerical dynamics show a re-entry even though the analytical bound is not satisfied, this primarily reflects conservativity of the bound, not evidence against the simulated mechanism.


\section{Corollary: Gain-Mediated Lucidity}

\begin{corollary}
\label{kor:lucidity}
If there exist a subset $S$ and an interval $I$ such that $\alpha_S(t)$ becomes small on $I$, the effective drives $\Lambda_i(t) = K\,G(t)\,g_i(t)$ for $i \in S$ do not fall below the reduced access threshold, and the reduced dynamics on $S$ remains in the access region with safety margin, then a transient re-entry of the full system onto $S$ is internally sufficiently possible within the model.
\end{corollary}

The mathematical core of the hypothesis:

\begin{quote}
\itshape
Within the GCC, a lucidity-like re-entry can be understood as finite-horizon persistence of a residual backbone regime under decreasing external perturbation load.
\end{quote}


\section{What This Appendix Delivers and What It Does Not}

\paragraph{Delivers.}
Well-posedness of the GCC as an SDE; controlled bounds for the core observables; finite-horizon stability; a genuine sufficient re-entry condition (Theorem~\ref{satz:reentry}).

\paragraph{Does not deliver.}
No mathematical proof that real terminal lucidity exists; no general theorem that every form of dementia shifts the access region; no proof of a universal threshold $K_c$; no identifiability result for real EEG/MEG data.


% ==================================================================
\section*{Supporting Information}
\addcontentsline{toc}{section}{Supporting Information}

\begin{sloppypar}
\noindent\textbf{S1 Code. GCC simulation and analysis code} (\texttt{gcc\_simulation\_v1.3.0.zip})\textbf{.} Complete Python codebase for the GCC dynamical model, per-subject EEG analysis pipeline, Chennu gamma baseline-holdout analysis (\texttt{run\_chennu\_baseline\_holdout\_from\_observables.py}), Hermann DoC proxy analysis (\texttt{run\_hermann2021\_gcc\_proxy\_analysis.py}), V3/V4 ensemble simulation (\texttt{test\_v2\_v3\_ensemble.py}), V6 retuning-control sweep (\texttt{test\_v5\_retuning\_control.py}), mixed-effects statistical analysis over the sweep (\texttt{test\_v5\_mixed\_effects.py}), and figure-generation scripts. Includes \texttt{requirements.txt} and \texttt{README.md} with reproduction instructions. Corresponds to release \texttt{v1.3.0} of the public repository \url{https://github.com/oktaytugan82/gcc_simulation}, prepared for Zenodo archival with DOI \href{https://doi.org/10.5281/zenodo.19744970}{10.5281/zenodo.19744970}.
\end{sloppypar}


% ==================================================================

\section*{Supplementary Clinical Plausibility Check: Hermann et al. 2021}
\addcontentsline{toc}{section}{Supplementary Clinical Plausibility Check: Hermann et al. 2021}

\noindent This supplementary analysis is included only as a clinical plausibility anchor. It is not a primary validation of the GCC and does not extract GCC observables from raw EEG or PET signals. The analysis uses the publicly released post-processed metadata table from \citet{hermann2021}, which contains diagnostic labels, six-month outcomes, PET-derived metabolic information, and EEG-derived classifier information. Because raw signal time series are not available in this table, the analysis can only test whether a simple combined PET-plus-EEG proxy behaves in the direction expected by the GCC's general claim that residual metabolic structure and functional signal evidence should be informative together.

\noindent For each participant, PET and EEG proxy variables were z-scored and averaged into a simple multimodal proxy. Receiver-operating-characteristic analyses were then computed for two clinically relevant endpoints: baseline MCS-/MCS+ versus VS/UWS diagnostic discrimination, and six-month MCS+/conscious versus other outcome discrimination. The resulting AUCs are shown in Supplementary Table~S1. We additionally used paired patient-level bootstrap resampling over complete cases to test whether the descriptive AUC advantages of the combined proxy were stable against paired sampling uncertainty.

\begin{table}[h]
\centering
\caption{Supplementary Table S1. Hermann et al.\ (2021) post-processed proxy analysis. AUC values are descriptive and based on the public metadata table only.}
\label{tab:hermann-proxy}
\small
\begin{tabular}{lccc}
\hline
Endpoint & PET proxy & EEG proxy & PET-plus-EEG proxy \\
\hline
Baseline MCS-/MCS+ vs.\ VS/UWS & 0.823 & 0.771 & 0.834 \\
Six-month MCS+/conscious vs.\ other & 0.720 & 0.690 & 0.762 \\
\hline
\end{tabular}
\end{table}

\noindent The combined proxy descriptively exceeds either unimodal proxy for both endpoints, but paired AUC-delta tests do not establish significant superiority. For baseline MCS-/MCS+ versus VS/UWS, $\Delta$AUC was $+0.017$ versus PET (95\% CI $-0.062$ to $+0.098$, $p = 0.67$) and $+0.063$ versus EEG (CI $-0.002$ to $+0.137$, $p = 0.059$). For six-month MCS+/conscious versus other outcome, $\Delta$AUC was $+0.040$ versus PET (CI $-0.054$ to $+0.140$, $p = 0.41$) and $+0.072$ versus EEG (CI $-0.007$ to $+0.160$, $p = 0.078$). This pattern is compatible with, but does not prove, the GCC's broader expectation that residual structural-metabolic capacity and functional signal organization are complementary in structurally degraded cohorts. The analysis should therefore be read as a supplementary bridge to the disorders-of-consciousness literature, not as evidence that the GCC access index is already a validated clinical biomarker.

\begin{figure}[h]
\centering
\includegraphics[width=0.92\linewidth]{data/hermann_2021_fdg_pet_eeg/hermann2021_gcc_proxy_diagnostic.png}
\caption{Supplementary Figure S5. Diagnostic Hermann et al.\ (2021) proxy analysis. PET, EEG, and PET-plus-EEG proxy scores are compared for baseline MCS-/MCS+ versus VS/UWS discrimination.}
\end{figure}

\begin{figure}[h]
\centering
\includegraphics[width=0.92\linewidth]{data/hermann_2021_fdg_pet_eeg/hermann2021_gcc_proxy_outcome.png}
\caption{Supplementary Figure S6. Prognostic Hermann et al.\ (2021) proxy analysis. PET, EEG, and PET-plus-EEG proxy scores are compared for six-month MCS+/conscious versus other outcome discrimination.}
\end{figure}

\begin{sloppypar}
\noindent Reproduction files for this supplementary check are stored under \path{data/hermann_2021_fdg_pet_eeg/}, including the public metadata table, the analysis script, derived features, summary JSON file, markdown report, and generated figures.
\end{sloppypar}

% ==================================================================

\section*{Supplementary Baseline-Window Holdout}
\addcontentsline{toc}{section}{Supplementary Baseline-Window Holdout}

\noindent The Cambridge raw-data bundle for \citet{chennu2016} was downloaded and its \texttt{datainfo.mat} table was used to verify state labels. For the theory-driven gamma-band pilot, each baseline recording was split into first-half calibration and second-half holdout. Regime bounds were estimated only from the first half of baseline and then applied to the held-out baseline half and to all sedation states. Results are summarized in Supplementary Table~S2.

\begin{table}[h]
\centering
\caption{Supplementary Table S2. Gamma-band baseline-window holdout for the Chennu et al.\ (2016) pilot.}
\label{tab:chennu-holdout}
\small
\begin{tabular}{lcccc}
\hline
Level & N & Mean $\Pi$ & SD & Median $\Pi$ \\
\hline
Baseline holdout & 20 & 0.658 & 0.032 & 0.659 \\
Mild & 20 & 0.599 & 0.085 & 0.626 \\
Moderate & 20 & 0.530 & 0.116 & 0.543 \\
Recovery & 20 & 0.572 & 0.153 & 0.648 \\
\hline
\end{tabular}
\end{table}

\noindent Paired tests against the baseline holdout remained significant for mild sedation (mean difference $+0.059$, Wilcoxon $p = 5.86 \times 10^{-4}$, $d_z = 0.75$, 16/20 declines), moderate sedation (mean difference $+0.128$, $p = 1.05 \times 10^{-4}$, $d_z = 1.05$, 17/20 declines), and recovery (mean difference $+0.086$, $p = 0.048$, $d_z = 0.55$, 13/20 declines). This is an internal within-recording holdout, not an independent baseline session.

% ==================================================================

\section*{Supplementary Figures}

\begin{figure}[h]
\centering
\includegraphics[width=0.98\linewidth]{pilot_trajectories_gamma.png}\\[0.3em]
\includegraphics[width=0.98\linewidth]{pilot_trajectories_alpha.png}
\caption{Supplementary Figure S1. Per-participant $\Pi$ trajectories across the four sedation levels. Top: gamma band. Bottom: alpha band. The dashed line at $\Pi = 0.5$ serves as orientation.}
\end{figure}

\begin{figure}[h]
\centering
\includegraphics[width=0.95\linewidth]{drowsy_responsive_diagnostic.png}
\caption{Supplementary Figure S2. Diagnostics for the drowsy--responsive question. Top: histograms of $\Pi$-declines (baseline minus moderate) per band. Bottom left: cross-band scatter of declines per participant; no discernible linear or cluster structure. Bottom right: ordered $\Pi$-declines per participant (sorted by gamma).}
\end{figure}

\begin{figure}[h]
\centering
\includegraphics[width=0.95\linewidth]{pilot_robustness.png}
\caption{Supplementary Figure S3. Robustness diagnostics for the pilot finding. Top: bootstrap distributions (10{,}000 resamples) of the mean $\Pi$-decline, gamma band (left) and alpha band (right). Bottom: leave-one-out $p$-values (Wilcoxon on the remaining 19 participants) as $-\log_{10}(p)$; both bands remain significant upon exclusion of any individual participant.}
\end{figure}

\begin{figure}[h]
\centering
\includegraphics[width=0.98\linewidth]{sweep_heatmap_gamma.png}\\[0.3em]
\includegraphics[width=0.98\linewidth]{sweep_heatmap_alpha.png}
\caption{Supplementary Figure S4. Sensitivity sweep over calibration hyperparameters. Heatmaps of $-\log_{10}(p)$ for the baseline-vs-moderate Wilcoxon test, per $\alpha$ quantile value (panels), $\tau_D$ (rows), and $\tau_M$ (x-axis within each panel). Top: gamma band. Bottom: alpha band. All 72 combinations remain significant; the direction of the effect is never reversed.}
\end{figure}

\begin{thebibliography}{1}

\bibitem[Chennu et~al.(2016)]{chennu2016}
Chennu,~S., O'Connor,~S., Adapa,~R., Menon,~D.\,K., \& Bekinschtein,~T.\,A.\ (2016).
Brain connectivity dissociates responsiveness from drug exposure during propofol-induced transitions of consciousness.
\textit{PLoS Computational Biology}, \textbf{12}(1), e1004669.

\bibitem[Hermann et~al.(2021)]{hermann2021}
Hermann,~B., Stender,~J., Habert,~M.-O., et~al.\ (2021).
Multimodal FDG-PET and EEG assessment improves diagnosis and prognostication of disorders of consciousness.
\textit{NeuroImage: Clinical}, \textbf{30}, 102601.
\url{https://doi.org/10.1016/j.nicl.2021.102601}

\end{thebibliography}

\end{document}
